\documentclass[12pt]{article}
\usepackage[margin=1in]{geometry}
\usepackage{enumitem}
\usepackage{titlesec}
\usepackage{parskip}
\usepackage{xcolor}
\usepackage{hyperref}
\usepackage{fontenc}

\title{\textbf{Cook and Levi 1990 Claude Guide}\\
\large Cook, Karen Schweers, and Margaret Levi, eds. \textit{The Limits of Rationality}.\\
Chicago: University of Chicago Press, 1990.}
\author{}
\date{}

\begin{document}

\maketitle
\tableofcontents
\newpage

%--------------------------------------------------
\section{Central Claims of the Volume}
%--------------------------------------------------

\begin{itemize}[leftmargin=*, itemsep=4pt]
    \item The volume is an interrogation of rational choice theory conducted by
    people largely sympathetic to it. The question is not whether rational
    choice is worth doing, but \textbf{where its boundary conditions lie}, and
    whether the failures at that boundary are patterned enough to be
    theorized.
    \item The organizing distinction, set up in Elster's opening chapter and
    running through the book, is between two very different kinds of failure:
    \begin{enumerate}
        \item \textbf{Indeterminacy} --- the theory fails to pick out a unique
        action, so it underdetermines behavior. This is the deeper problem,
        because it is a defect in the theory rather than in the agent.
        \item \textbf{Irrationality} --- agents fail to do what the theory
        prescribes. Damaging, but survivable, and empirically tractable.
    \end{enumerate}
    \item A second organizing claim is that the rational actor model is
    \textbf{incomplete rather than wrong}: it takes preferences as given and
    says nothing about where they come from, so it must be supplemented by
    accounts of \textbf{preference formation, norms, and institutions}. Parts II
    and III supply those supplements.
    \item There is a recurring \textbf{level-of-analysis} argument: individual
    choice may be demonstrably non-rational while aggregate outcomes still
    match the competitive predictions of the model (Plott), which means
    evidence of individual irrationality does not automatically defeat rational
    choice as a tool for institutional or market-level explanation.
    \item The institutional chapters converge on the view that
    \textbf{institutions do the work that incentives alone cannot} --- solving
    collective action problems, generating and sustaining norms, and making
    cooperation self-enforcing. Rationality is thus \emph{socially embedded}
    rather than a property of isolated agents.
\end{itemize}

%--------------------------------------------------
\section{Organization}
%--------------------------------------------------

Most chapters are paired with a short critical \textbf{comment} by another
contributor, so the book reads as a structured debate rather than a
collection. Editors' introduction is by Levi, Cook, Jodi A. O'Brien, and
Howard Faye. Three parts: the theory itself, preference formation and norms,
and institutions.

%--------------------------------------------------
\section{Part I: The Theory of Rational Choice}
%--------------------------------------------------

\subsection*{1. Jon Elster, ``When Rationality Fails'' (comment: Geoffrey Brennan)}
\begin{itemize}[leftmargin=*, itemsep=3pt]
    \item Sets the volume's agenda by separating \textbf{indeterminacy} from
    \textbf{irrationality}, and argues indeterminacy is the more serious
    charge: where the theory yields multiple or no prescriptions, it cannot be
    rescued by better data about agents.
    \item Key illustration: the \textbf{infinite regress in information
    gathering}. Deciding how much information to collect before choosing is
    itself a decision requiring information, and the regress has no rational
    stopping point.
    \item Strategic settings with \textbf{multiple equilibria} are the other
    standing source of indeterminacy --- game theory frequently cannot say
    which equilibrium obtains.
    \item Brennan's comment presses on what exactly we should want rationality
    to deliver, and whether the demands Elster places on it are ones any
    theory could meet.
\end{itemize}

\subsection*{2. Amos Tversky and Daniel Kahneman, ``Rational Choice and the Framing of Decisions''}
\begin{itemize}[leftmargin=*, itemsep=3pt]
    \item The most-cited chapter. Choices systematically violate
    \textbf{invariance} (logically equivalent descriptions should yield the same
    choice) and sometimes \textbf{dominance} --- and these are not random
    error but predictable, reversible effects of \textbf{framing}.
    \item Prospect theory's apparatus: outcomes are evaluated as
    \textbf{gains and losses from a reference point} rather than as final
    states; \textbf{loss aversion}; and the \textbf{reflection effect} ---
    risk aversion over gains, risk seeking over losses.
    \item Cases: the \textbf{Asian disease problem} (identical mortality
    outcomes framed as lives saved vs.\ lives lost reverse the majority
    preference); \textbf{mental accounting} examples such as the theater
    ticket and the jacket-and-calculator problems, where the same sum is
    treated differently depending on the account it is booked to.
    \item Directly threatens the descriptive adequacy of expected utility while
    leaving its normative status open.
\end{itemize}

\subsection*{3. Mark J. Machina, ``Choice Under Uncertainty: Problems Solved and Unsolved'' (comment: John Broome)}
\begin{itemize}[leftmargin=*, itemsep=3pt]
    \item A survey of the non-expected-utility research program: the
    documented anomalies do not require abandoning formal choice theory, only
    weakening the \textbf{independence axiom}.
    \item Cases: the \textbf{Allais paradox} (systematic independence
    violations) and the \textbf{Ellsberg paradox} (ambiguity aversion ---
    preference for known over unknown probabilities).
    \item Machina's ``generalized expected utility'' preserves smooth
    preferences and most analytic results while accommodating the anomalies.
    The constructive counterweight to Tversky and Kahneman.
    \item Broome's comment shifts to the normative question of whether a
    rational agent \emph{ought} to maximize expected utility.
\end{itemize}

\subsection*{4. Charles R. Plott, ``Rational Choice in Experimental Markets'' (comment: Cook and O'Brien)}
\begin{itemize}[leftmargin=*, itemsep=3pt]
    \item The strongest defense in the volume. In \textbf{experimental
    double-auction markets}, prices and quantities converge on competitive
    equilibrium predictions rapidly and reliably --- even with few traders,
    incomplete information, and individuals who are demonstrably not
    maximizers.
    \item Implication: \textbf{institutions aggregate}. Market structure can
    produce rational outcomes out of imperfectly rational agents, so
    individual-level anomalies need not propagate upward.
    \item Cook and O'Brien's comment is the volume's pivot: market-level
    predictive success does not license claims about individual decision
    processes, and the two levels should not be conflated.
\end{itemize}

%--------------------------------------------------
\section{Part II: Preference Formation and the Role of Norms}
%--------------------------------------------------

\subsection*{5. George J. Stigler and Gary S. Becker, ``De Gustibus Non Est Disputandum'' (comment: Robert E. Goodin)}
\begin{itemize}[leftmargin=*, itemsep=3pt]
    \item The hard-line position, reprinted as the foil for the rest of the
    part: \textbf{tastes are stable over time and broadly similar across
    people}. Explain behavior with prices, income, and constraints --- never by
    positing a change in preferences, which they treat as an admission of
    explanatory defeat.
    \item Mechanism: \textbf{household production} and accumulated
    \textbf{consumption capital}. Apparent taste change is really a change in
    the shadow price of a commodity the household produces.
    \item Cases: \textbf{music appreciation} (``beneficial addiction'' --- past
    listening lowers the cost of present enjoyment), \textbf{advertising},
    \textbf{fashion}, and \textbf{custom and tradition}, each re-described
    without invoking shifting tastes.
    \item Goodin's ``De Gustibus Non Est Explanandum'' is the direct rebuttal:
    the strategy makes preferences unexplainable by fiat rather than
    explaining them.
\end{itemize}

\subsection*{6. Michael Taylor, ``Cooperation and Rationality'' (comment: Michael Hechter)}
\begin{itemize}[leftmargin=*, itemsep=3pt]
    \item Revisits the collective action problem: the pessimism of the one-shot
    Prisoner's Dilemma dissolves under \textbf{repeated play} among a stable
    population, where \textbf{conditional cooperation} sustained by the shadow
    of the future can be an equilibrium.
    \item Conditions favoring cooperation: \textbf{small groups}, ongoing
    interaction, low discount rates, and \textbf{community} --- shared beliefs
    and dense direct relations. Anticipates Ostrom's design principles.
    \item Hechter's comment is the skeptical counterweight: supergame results
    are too permissive (folk-theorem multiplicity) to explain real-world
    collective action, and the indeterminacy Elster identified resurfaces here.
\end{itemize}

\subsection*{7. James S. Coleman, ``Norm-Generating Structures'' (comment: Stephen J. Majeski)}
\begin{itemize}[leftmargin=*, itemsep=3pt]
    \item Treats norms as things to be \emph{explained}, not assumed. A
    \textbf{demand for a norm} arises when actions impose \textbf{externalities}
    that cannot be resolved by direct bargaining.
    \item Demand is not sufficient: realizing a norm requires \textbf{social
    structure that makes sanctioning feasible}, since sanctioning is itself a
    second-order collective action problem. Coleman's answer is
    \textbf{closure} --- dense, interconnected relations letting beneficiaries
    coordinate and share sanctioning costs.
    \item Keeps rational choice microfoundations while generating genuinely
    social outcomes; connects directly to his social capital work.
    \item Majeski proposes an alternative route to norm generation and
    maintenance.
\end{itemize}

%--------------------------------------------------
\section{Part III: Institutions}
%--------------------------------------------------

\subsection*{8. Arthur S. Stinchcombe, ``Reason and Rationality'' (comment: Andrew Abbott)}
\begin{itemize}[leftmargin=*, itemsep=3pt]
    \item Distinguishes \textbf{rationality} (maximizing given ends) from
    \textbf{reason} --- deliberation, argument, and the giving and assessing of
    reasons. Much institutional life runs on the latter.
    \item Institutions such as courts, scientific disciplines, and
    bureaucracies are organized to produce and evaluate reasons, and a
    maximizing model of the individual misses what they do.
    \item Abbott's comment engages the distinction from a sociological
    standpoint.
\end{itemize}

\subsection*{9. Gary J. Miller, ``Managerial Dilemmas'' (comment: Robert H. Bates and William T. Bianco)}
\begin{itemize}[leftmargin=*, itemsep=3pt]
    \item Hierarchies cannot fully solve the \textbf{team production} problem
    with incentives alone: when output is joint and individual contribution is
    unobservable, no budget-balancing incentive scheme achieves efficiency
    (the Holmstr\"om result).
    \item Because the problem is formally unsolvable by contract,
    \textbf{leadership, credible commitment, and organizational culture} do
    real work --- an argument that political and normative factors are
    \emph{required}, not residual.
    \item Bates and Bianco extend the leadership argument within a rational
    choice frame.
\end{itemize}

\subsection*{10. Russell Hardin, ``The Social Evolution of Cooperation'' (comment: Carol A. Heimer)}
\begin{itemize}[leftmargin=*, itemsep=3pt]
    \item Many cooperation problems are better modeled as \textbf{coordination}
    than as Prisoner's Dilemmas. Once a convention is established it is
    largely \textbf{self-enforcing}, since no one gains by unilateral
    deviation.
    \item Cooperative norms and conventions therefore \textbf{evolve} through
    ongoing interaction rather than requiring design or external enforcement.
    \item Heimer's comment presses on the limits of the evolutionary account.
\end{itemize}

\subsection*{11. Douglass C. North, ``Institutions and Their Consequences for Economic Performance''}
\begin{itemize}[leftmargin=*, itemsep=3pt]
    \item Institutions are the \textbf{rules of the game} --- formal rules,
    informal constraints, and enforcement characteristics --- and they exist
    because \textbf{transaction costs} are positive.
    \item The central puzzle is the persistence of \textbf{inefficient
    institutions}: because returns are shaped by existing rules and by the
    organizations those rules create, \textbf{path dependence} can lock in
    arrangements that do not maximize aggregate wealth.
    \item Compact statement of the argument of \textit{Institutions,
    Institutional Change and Economic Performance}, also 1990 --- read
    alongside your North notes.
\end{itemize}

\subsection*{12. Margaret Levi, ``A Logic of Institutional Change''}
\begin{itemize}[leftmargin=*, itemsep=3pt]
    \item The editors' closing statement. Institutional change is driven by
    rulers acting as \textbf{revenue maximizers} subject to constraints:
    relative bargaining power, transaction costs, and discount rates.
    \item \textbf{Quasi-voluntary compliance} is the key concept: compliance is
    voluntary in that it is not usually coerced, but quasi- in that
    non-compliance is punishable. Rulers prefer it because \textbf{enforcement is
    expensive}, which gives them reason to make credible commitments and
    supply something in return.
    \item Compliance is \textbf{contingent} --- it depends on rulers keeping
    their bargains and on citizens believing others are also complying, so
    normative and strategic considerations are entangled rather than opposed.
    \item Cases drawn from her work on \textbf{taxation and revenue extraction}
    (see \textit{Of Rule and Revenue}, 1988), and later extended to
    \textbf{conscription}.
\end{itemize}

%--------------------------------------------------
\section{Connections to Your Other Readings}
%--------------------------------------------------

\begin{itemize}[leftmargin=*, itemsep=4pt]
    \item \textbf{Olson 1965}: the target of Part II. Taylor's repeated-game
    argument and Coleman's norm-generating structures are both direct responses
    to Olsonian pessimism about collective action.
    \item \textbf{Ostrom 1990} (same year): Taylor's conditions for cooperation
    --- small stable groups, repeated interaction, community --- are close to
    Ostrom's design principles, arrived at from theory rather than fieldwork.
    The strongest single pairing on this list.
    \item \textbf{North 1990}: chapter 11 \emph{is} North's argument in
    miniature; useful for citing him on transaction costs and path dependence
    without leaving this volume.
    \item \textbf{Putnam 1993}: Coleman on closure is the direct theoretical
    ancestor of Putnam's social capital.
    \item \textbf{Scott 1985}: peasant behavior that is unintelligible under a
    narrow maximizing model becomes intelligible once norms and dignity enter,
    which is Part II's point made ethnographically.
    \item \textbf{Svolik 2012}: authoritarian power-sharing turns on credible
    commitment under repeated interaction --- Miller's and Hardin's logic
    applied to dictatorship.
    \item \textbf{Weber}: the \textit{Zweckrational}/\textit{Wertrational}
    distinction is the ancestor of the whole debate, and Stinchcombe's
    reason/rationality split is its clearest modern descendant.
    \item \textbf{Huber and Shipan 2002}: Miller's managerial dilemmas are the
    delegation and discretion problem in the firm rather than the legislature.
\end{itemize}

%--------------------------------------------------
\section{Likely Exam Uses}
%--------------------------------------------------

\begin{itemize}[leftmargin=*, itemsep=4pt]
    \item The \textbf{indeterminacy/irrationality distinction} (Elster) is the
    single most portable idea here --- it lets you criticize a rational choice
    argument precisely rather than gesturally.
    \item \textbf{Plott vs.\ Cook and O'Brien} is the cleanest citation for
    level-of-analysis problems: aggregate predictive success versus individual
    process claims.
    \item \textbf{Levi on quasi-voluntary compliance} is directly usable in any
    question about state capacity, taxation, or legitimacy, and connects the
    coercion and consent literatures.
    \item As a \textbf{bridge}: the volume lets you move from rational choice to
    institutional or normative explanation without the shift looking like a
    change of subject.
\end{itemize}

\end{document}
